\chapter{Αποτελέσματα}

\section{Screened Εκτιμήσεις}

Με βάση τα κριτήρια screening της Ενότητας~\ref{amb:preprocessing}, το πρώτο sweep τάξεων, $n \in \{2,4,\dots,30\}$, αποδίδει 1397 εκτιμήσεις, ενώ το δεύτερο sweep τάξεων, $n \in \{10,15,\dots,45\}$, αποδίδει 1391 εκτιμήσεις. Η κατανομή τους ανά περιοχή ελέγχου παρουσιάζεται στον Πίνακα~\ref{tab:amb:screened_estimates}.

\begin{table}[htbp]
    \centering
    \caption{Πλήθος screened εκτιμήσεων N4SID ανά order sweep και περιοχή ελέγχου.}
    \label{tab:amb:screened_estimates}
    \begin{tabular}{ccccc}
        \toprule
        \textbf{Order sweep} & \textbf{Περιοχή 1} & \textbf{Περιοχή 2} & \textbf{Περιοχή 3} & \textbf{Σύνολο} \\
        \midrule
        $n \in \{2,4,\dots,30\}$ & 565 & 280 & 552 & 1397 \\
        $n \in \{10,15,\dots,45\}$ & 548 & 272 & 571 & 1391 \\
        \bottomrule
    \end{tabular}
\end{table}

Το Σχήμα~\ref{fig:amb:screened_modal_grid} απεικονίζει τις screened εκτιμήσεις στο επίπεδο $(\sigma,f)$, ανά order sweep και περιοχή ελέγχου. Τα κοινά όρια των αξόνων επιτρέπουν την άμεση οπτική σύγκριση των modal συγκεντρώσεων, ενώ οι δείκτες αναφοράς αντιστοιχούν στους ρυθμούς του \textit{PowerFactory} του Πίνακα~\ref{tab:amb:pf_reference_modes}.

\clearpage
\begin{figure}[p]
    \centering
    \includegraphics[width=\textwidth,height=0.86\textheight,keepaspectratio]{\ambientfigurepath Ambient_Mag0.1_T600s_dt10ms_seed1997/plots/pdf/screened_modal_maps_3x2_grid.pdf}
    \caption{Screened εκτιμήσεις N4SID στο επίπεδο $(\sigma,f)$ ανά order sweep και περιοχή ελέγχου. Οι χρωματισμοί διακρίνουν τα σήματα τάσης και ρεύματος, ενώ οι λευκοί ρόμβοι δηλώνουν τους ρυθμούς αναφοράς του \textit{PowerFactory} του Πίνακα~\ref{tab:amb:pf_reference_modes}.}
    \label{fig:amb:screened_modal_grid}
\end{figure}
\clearpage

Είναι εμφανές ότι τα δύο sweeps εμφανίζουν παρόμοια κατανομή των εκτιμήσεων και στις τρεις περιοχές ελέγχου, καθώς και ότι οι εκτιμήσεις του N4SID συγκεντρώνονται κοντά στις συχνότητες των ρυθμών αναφοράς του \textit{PowerFactory}, επομένως η μέθοδος αποδίδει ικανοποιητικά τη θέση των ηλεκτρομηχανικών ρυθμών ως προς τη συχνότητα $f$. Αντίθετα, παρατηρείται μεγαλύτερη διασπορά κατά τον άξονα της απόσβεσης $\sigma$. Η διασπορά αυτή αναδεικνύει ότι η εκτίμηση της απόσβεσης είναι λιγότερο σταθερή από την εκτίμηση της συχνότητας και αιτιολογεί την επόμενη εφαρμογή συσταδοποίησης.


\section{Αποτελέσματα Συσταδοποίησης ανά Μέθοδο}
\label{sec:amb:clustering-by-method}

Η παρούσα ενότητα παρουσιάζει, για κάθε αλγόριθμο, τις ρυθμίσεις που
επιλέχθηκαν από τη διαδικασία tuning της Ενότητας~\ref{amb:clustering-tuning}.
Οι πίνακες συνοψίζουν τις επιλεγμένες ρυθμίσεις για κάθε order sweep και
περιοχή ελέγχου. Η ερμηνεία των αποτελεσμάτων βασίζεται στο πλήθος των
συστάδων, στον χειρισμό του θορύβου και στο \textit{silhouette score}, ενώ η
άμεση σύγκριση όλων των μεθόδων ακολουθεί στην επόμενη ενότητα.

\subsection{K-means}

Για τον \textit{k-Means}, η επιλογή του αριθμού συστάδων βασίστηκε στη
μέγιστη τιμή silhouette και στο φιλτράρισμα της Ενότητας~\ref{amb:silhouette-filtering}. Οι επιλεγμένες λύσεις περιλαμβάνουν από δύο έως
πέντε συστάδες και όλα τα σημεία αποδίδονται σε κάποια συστάδα καθώς ο αλγόριθμος δεν κάνει διαχωρισμό των σημείων σε θόρυβο αλλά και καμία συστάδα δεν προέκυψε με μέσο silhoutte score μικρότερο του 0.25. Όλες οι τιμές των Silhouette Score είναι παρόμοιες, γεγονός που υποδηλώνει παρόμοια συνοχή των συστάδων σε όλες τις περιοχές ελέγχου με τις τιμές στην Περιοχή 2 να είναι ελαφρώς υψηλότερες και στα δύο sweeps. Οι επιλεγμένες τιμές του \textit{k} παρουσιάζονται στον Πίνακα~\ref{tab:amb:kmeans-tuning}.

\begin{table}[htbp]
    \centering
    \caption{Επιλεγμένες ρυθμίσεις \textit{k-Means}.}
    \label{tab:amb:kmeans-tuning}
    \begin{tabular}{cccc}
        \toprule
        Sweep & Περιοχή & $k$ & Silhouette Score \\
        \midrule
        orders1 & 1 & 5 & 0.640 \\
        orders1 & 2 & 4 & 0.700 \\
        orders1 & 3 & 4 & 0.650 \\
        orders2 & 1 & 2 & 0.609 \\
        orders2 & 2 & 5 & 0.681 \\
        orders2 & 3 & 2 & 0.621 \\
        \bottomrule
    \end{tabular}
\end{table}

\subsection{K-medoids}

Ο \textit{k-Medoids} ρυθμίστηκε με το ίδιο εύρος τιμών $K$ και το ίδιο
κριτήριο επιλογής με τον \textit{k-Means} όπως περιγράφηκε στην Ενότητα~\ref{amb:clustering-tuning}. Παρατηρείται ότι οι περισσότερες επιλεγμένες τιμές του $k$ είναι ίδιες με αυτές του \textit{k-Means}, ενώ εξαιρετική ομοιότητα παρατηρείται και στις τιμές του silhouette score. Όμοια με τον k-Means, δεν προέκυψε καμία συστάδα με Silhouette Score μικρότερο του 0.25 οπότε δεν χαρακτηρίζεται καμία εκτίμηση ως θόρυβος. Οι επιλεγμένες τιμές του \textit{k} παρουσιάζονται στον Πίνακα~\ref{tab:amb:kmedoids-tuning}.

\begin{table}[htbp]
    \centering
    \caption{Επιλεγμένες ρυθμίσεις \textit{k-Medoids}.}
    \label{tab:amb:kmedoids-tuning}
    \begin{tabular}{cccc}
        \toprule
        Sweep & Περιοχή & $k$ & Silhouette Score \\
        \midrule
        orders1 & 1 & 5 & 0.634 \\
        orders1 & 2 & 4 & 0.702 \\
        orders1 & 3 & 4 & 0.654 \\
        orders2 & 1 & 2 & 0.609 \\
        orders2 & 2 & 4 & 0.664 \\
        orders2 & 3 & 4 & 0.621 \\
        \bottomrule
    \end{tabular}
\end{table}

\subsection{DBSCAN}

Ο DBSCAN επέλεξε σε όλες τις περιπτώσεις, εκτός μίας, την ελάχιστη τιμή της παραμέτρου $p_e$ ($p_e = 0.050$). Ωστόσο, δεν κρίνεται σκόπιμο να μειωθεί το κατώφλι $p_e$, καθώς θα αυξανόταν περαιτέρω το ποσοστό των εκτιμήσεων που χαρακτηρίζονται ως θόρυβος. Ακόμα, η μέθοδος παρήγαγε από δύο έως πέντε συστάδες σε κάθε περίπτωση με υψηλά silhouette scores, γεγονός που υποδηλώνει πολύ συνεκτικές συστάδες, με μειονέκτημα όμως την απόρριψη μεγάλου μέρους των εκτιμήσεων ως θόρυβο. Οι επιλεγμένες τιμές των υπερπαραμέτρων παρουσιάζονται στον Πίνακα~\ref{tab:amb:dbscan-tuning}.

\begin{table}[htbp]
    \centering
    \caption{Επιλεγμένες ρυθμίσεις DBSCAN.}
    \label{tab:amb:dbscan-tuning}
    \begin{tabular}{cccccccc}
        \toprule
        Sweep & Περιοχή & $p_e$ & $p_m$ & $N_{\mathrm{pts}}$ & Clusters & Noise [\%] & Silhouette \\
        \midrule
        orders1 & 1 & 0.050 & 0.40 & 48 & 2 & 59.3 & 0.876 \\
        orders1 & 2 & 0.050 & 0.40 & 24 & 3 & 50.0 & 0.927 \\
        orders1 & 3 & 0.050 & 0.35 & 42 & 2 & 70.5 & 0.936 \\
        orders2 & 1 & 0.050 & 0.40 & 26 & 5 & 39.1 & 0.716 \\
        orders2 & 2 & 0.060 & 0.35 & 12 & 4 & 21.3 & 0.774 \\
        orders2 & 3 & 0.050 & 0.35 & 23 & 5 & 35.2 & 0.752 \\
        \bottomrule
    \end{tabular}
\end{table}

\subsection{OPTICS}

Ο αλγόριθμος OPTICS επέλεξε από δύο έως επτά συστάδες στις εξεταζόμενες περιπτώσεις, με επίσης υψηλά silhouette scores. Όμως, όπως
και στον DBSCAN, οι υψηλές τιμές silhouette συνοδεύονται από μεγάλο ποσοστό
σημείων θορύβου σε αρκετές περιπτώσεις. Επομένως, οι ρυθμίσεις με τη μεγαλύτερη τιμή silhouette
περιγράφουν πολύ συνεκτικές ομάδες εκτιμήσεων, χωρίς όμως να καλύπτουν το
σύνολο του διαθέσιμου δείγματος. Οι επιλεγμένες τιμές των υπερπαραμέτρων παρουσιάζονται στον Πίνακα~\ref{tab:amb:optics-tuning}.

\begin{table}[htbp]
    \centering
    \caption{Επιλεγμένες ρυθμίσεις OPTICS.}
    \label{tab:amb:optics-tuning}
    \begin{tabular}{cccccccc}
        \toprule
        Sweep & Περιοχή & $p_m$ & \texttt{xi} & \texttt{min\_samples} & Clusters & Noise [\%] & Silhouette \\
        \midrule
        orders1 & 1 & 0.20 & 0.18 & 24 & 2 & 87.1 & 0.952 \\
        orders1 & 2 & 0.30 & 0.18 & 18 & 5 & 44.3 & 0.922 \\
        orders1 & 3 & 0.15 & 0.13 & 18 & 4 & 62.9 & 0.935 \\
        orders2 & 1 & 0.25 & 0.17 & 16 & 6 & 63.7 & 0.857 \\
        orders2 & 2 & 0.35 & 0.14 & 12 & 7 & 41.5 & 0.817 \\
        orders2 & 3 & 0.35 & 0.13 & 23 & 4 & 73.9 & 0.952 \\
        \bottomrule
    \end{tabular}
\end{table}

\subsection{HDBSCAN}

Ο HDBSCAN διατηρεί πολύ μεγαλύτερο μέρος των εκτιμήσεων σε σχέση με τις δύο
προηγούμενες density-based μεθόδους, επιλέγοντας όμως από τέσσερις έως εννέα
συστάδες. Οι τιμές silhouette παραμένουν υψηλές σε όλες τις περιοχές. Oπτικά η μέθοδος φαίνεται να ανακτά όλες τις βασικές διακριτές συγκεντρώσεις των εκτιμήσεων στο επίπεδο $(\sigma,f)$ ενώ παράλληλα απορρίπτει ως θόρυβο τις εκτιμήσεις που δεν ανήκουν σε πυκνές περιοχές. Η μέθοδος \texttt{eom} επιλέχθηκε σε όλες τις περιπτώσεις, με μοναδική εξαίρεση το δεύτερο sweep της Περιοχής 3, όπου επιλέχθηκε η πολιτική \texttt{leaf}. Οι επιλεγμένες τιμές των υπερπαραμέτρων παρουσιάζονται στον Πίνακα~\ref{tab:amb:hdbscan-tuning}.

\begin{table}[htbp]
    \centering
    \caption{Επιλεγμένες ρυθμίσεις HDBSCAN.}
    \label{tab:amb:hdbscan-tuning}
    \resizebox{\textwidth}{!}{%
    \begin{tabular}{ccccccccc}
        \toprule
        Sweep & Περιοχή & $p_e$ & $p_m$ & \texttt{min\_cluster\_size} & Method & Clusters & Noise [\%] & Silhouette \\
        \midrule
        orders1 & 1 & 0.010 & 0.35 & 42 & eom & 4 & 29.2 & 0.796 \\
        orders1 & 2 & 0.010 & 0.20 & 12 & eom & 7 & 21.4 & 0.782 \\
        orders1 & 3 & 0.010 & 0.20 & 24 & eom & 5 & 25.4 & 0.816 \\
        orders2 & 1 & 0.010 & 0.30 & 20 & eom & 9 & 34.1 & 0.741 \\
        orders2 & 2 & 0.065 & 0.40 & 13 & eom & 4 & 12.5 & 0.749 \\
        orders2 & 3 & 0.010 & 0.40 & 26 & leaf & 6 & 39.2 & 0.742 \\
        \bottomrule
    \end{tabular}
    }
\end{table}

\subsection{Gaussian Mixture Models}

Για τον Gaussian Mixture Models, οι υπερπαράμετροι επιλέχθηκαν με κριτήριο την
ελαχιστοποίηση του BIC και όχι με βάση το silhouette Score όπως στις υπόλοιπες μεθόδους. Η αρχική λύση
επιλέγει εννέα ή δέκα Gaussian components/clusters σε κάθε περίπτωση και αποδίδει όλες τις εκτιμήσεις
σε κάποιο component/cluster. Στη συνέχεια όμως το φίλτρο της Ενότητας~\ref{amb:silhouette-filtering} απορρίπτει ως θόρυβο έως και πέντε από τις αρχικές συστάδες. Τελικά, οπτικά φαίνεται να  ανακτά όλες τις εμφανείς συγκεντρώσεις εκτιμήσεων παρόμοια με τον HDBSCAN. Οι αντίστοιχες τιμές Silhouette πριν και μετά το φιλτράρισμα παρουσιάζονται μόνο συμπληρωματικά και δεν αποτελούν κριτήριο επιλογής του μοντέλου. Οι επιλεγμένες τιμές των υπερπαραμέτρων, οι συστάδες και τα Silhouette Score παρουσιάζονται στον Πίνακα~\ref{tab:amb:gmm-tuning}.

\begin{table}[htbp]
    \centering
    \caption{Επιλεγμένες ρυθμίσεις Gaussian Mixture Models.}
    \label{tab:amb:gmm-tuning}
    \begin{tabular}{ccccccc}
        \toprule
        Sweep & Περιοχή & $k$ & BIC & \makecell{Silhouette\\(tuning)} & \makecell{Τελικές\\συστάδες} & \makecell{Τελικό\\silhouette}\\
        \midrule
        orders1 & 1 & 10 & 2465.84 & 0.363 & 7 & 0.621 \\
        orders1 & 2 & 10 & 959.29 & 0.454 & 7 & 0.734 \\
        orders1 & 3 & 10 & 2337.83 & 0.374 & 6 & 0.718 \\
        orders2 & 1 & 10 & 2340.35 & 0.083 & 6 & 0.705 \\
        orders2 & 2 & 10 & 1142.22 & 0.497 & 8 & 0.587 \\
        orders2 & 3 & 9 & 2559.48 & 0.143 & 4 & 0.821 \\
        \bottomrule
    \end{tabular}
\end{table}

\subsection{Agglomerative Hierarchical Clustering}

Η Agglomerative Hierarchical Clustering αρχικά επίσης αποδίδει όλες τις εκτιμήσεις σε
συστάδες, χωρίς χαρακτηρισμό σημείων ως θορύβου. Οι ρυθμίσεις που
μεγιστοποιούν το silhouette οδηγούν σε τρεις έως επτά συστάδες, σε δύο όμως περιπτώσεις το post-processing της Ενότητας~\ref{amb:silhouette-filtering} αφαιρεί συστάδες χαμηλής συνοχής. Οι επιλεγμένες τιμές των υπερπαραμέτρων καθώς και ο αρχικός και τελικός αριθμός συστάδων καθώς και τα αντίστοιχα Silhouette Score παρουσιάζονται στον Πίνακα~\ref{tab:amb:agglomerative-tuning}.

\begin{table}[htbp]
    \centering
    \caption{Επιλεγμένες ρυθμίσεις Agglomerative Hierarchical Clustering.}
    \label{tab:amb:agglomerative-tuning}
    \begin{tabular}{cccccccc}
        \toprule
        Sweep & Περιοχή & $p_e$ & Linkage & \makecell{Clusters\\(tuning)} & \makecell{Silhouette\\(tuning)} & \makecell{Τελικές\\συστάδες} & \makecell{Τελικό\\silhouette} \\
        \midrule
        orders1 & 1 & 0.36 & average & 7 & 0.626 & 5 & 0.633 \\
        orders1 & 2 & 0.54 & complete & 4 & 0.685 & 4 & 0.685 \\
        orders1 & 3 & 0.73 & complete & 6 & 0.633 & 6 & 0.633 \\
        orders2 & 1 & 0.79 & average & 3 & 0.569 & 2 & 0.587 \\
        orders2 & 2 & 0.49 & average & 4 & 0.670 & 4 & 0.670 \\
        orders2 & 3 & 0.67 & complete & 6 & 0.607 & 6 & 0.607 \\
        \bottomrule
    \end{tabular}
\end{table}


\section{Συγκριτικά Αποτελέσματα όλων των Μεθόδων Συσταδοποίησης}

Η διαδικασία tuning οδηγεί σε διαφορετικό πλήθος συστάδων, ανάλογα με τη
μέθοδο, την περιοχή ελέγχου και το order sweep. Το πλήθος αυτό αποτελεί
χρήσιμη πρώτη ένδειξη της συμπεριφοράς κάθε μεθόδου: πολύ μικρές τιμές
ενδέχεται να συγχωνεύουν ξεχωριστές modal συγκεντρώσεις, ενώ πολύ μεγάλες
τιμές μπορεί να υποδεικνύουν υπερκατατμημένη περιγραφή των εκτιμήσεων.

Ο Πίνακας~\ref{tab:amb:clustering-comparison} συγκεντρώνει το επιλεγμένο
πλήθος συστάδων ανά order sweep και περιοχή ελέγχου. Παρουσιάζεται επίσης
το πλήθος των ρυθμών αναφοράς του \textit{PowerFactory} που σχετίζονται με
κάθε περιοχή, σύμφωνα με τον Πίνακα~\ref{tab:amb:pf_reference_modes}. Οι
αναλυτικές ρυθμίσεις που οδήγησαν στις συγκεκριμένες επιλογές παρουσιάστηκαν
στην Ενότητα~\ref{sec:amb:clustering-by-method}.

\clearpage
\begin{table}[p]
    \centering
    \rotatebox{90}{%
    \begin{minipage}{0.92\textheight}
    \centering
    \captionof{table}{Πλήθος συστάδων της επιλεγμένης ρύθμισης ανά order sweep, περιοχή ελέγχου και μέθοδο.}
    \label{tab:amb:clustering-comparison}
    \resizebox{\linewidth}{!}{%
    \begin{tabular}{ccc|ccccccc}
        \toprule
        \textbf{Sweep} & \textbf{Περ.} & \shortstack{\textbf{Ref.}\\\textbf{modes}} & \shortstack{\textbf{k-}\\\textbf{Means}} & \shortstack{\textbf{k-}\\\textbf{Medoids}} & \textbf{DBSCAN} & \textbf{OPTICS} & \textbf{HDBSCAN} & \textbf{GMM} & \textbf{AHC} \\
        \midrule
        \multirow{3}{*}{orders1} & 1 & 5 & 5 & 5 & 2 & 2 & 4 & 7 & 5 \\
        & 2 & 3 & 4 & 4 & 3 & 5 & 7 & 7 & 4 \\
        & 3 & 6 & 4 & 4 & 2 & 4 & 5 & 6 & 6 \\
        \midrule
        \multirow{3}{*}{orders2} & 1 & 5 & 2 & 2 & 5 & 6 & 9 & 6 & 2 \\
        & 2 & 3 & 5 & 4 & 4 & 7 & 4 & 8 & 4 \\
        & 3 & 6 & 2 & 4 & 5 & 4 & 6 & 4 & 6 \\
        \bottomrule
    \end{tabular}
    }
    \end{minipage}
    }
\end{table}
\clearpage

Οι μέθοδοι k-Means, k-Medoids, AHC και GMM αρχικά αποδίδουν όλες τις
εκτιμήσεις σε συστάδες. Και οι τέσσερις αλγόριθμοι επιλέγουν από δύο έως οκτώ συστάδες, πλήθος που
είναι γενικά συγκρίσιμο με τους τρεις έως έξι ρυθμούς αναφοράς των επιμέρους
περιοχών. Ακόμα, οι DBSCAN και OPTICS επιλέγουν από δύο έως πέντε και από δύο έως επτά
συστάδες αντίστοιχα, αλλά χαρακτηρίζουν ως θόρυβο μεγάλο μέρος των εκτιμήσεων, πετυχαίνοντας έτσι σαφώς υψηλότερα silhouette scores.
Επομένως, οι υψηλές τιμές silhouette των δύο αυτών μεθόδων πρέπει να
ερμηνεύονται μαζί με το ποσοστό των αποδιδόμενων σημείων και δεν
αποτελούν από μόνες τους ένδειξη πληρέστερης αναπαράστασης των modal
συγκεντρώσεων. Τέλος, ο HDBSCAN αποτελεί ενδιάμεση περίπτωση. Διατηρεί σημαντικά μεγαλύτερο μέρος
του δείγματος από τις άλλες density-based μεθόδους, αλλά εντοπίζει περισσότερες συστάδες (τέσσερις
έως εννέα). Η τελική αξιολόγηση απαιτεί επομένως συνδυασμό του πλήθους συστάδων, της κάλυψης των
εκτιμήσεων και της οπτικής αντιστοίχισης με τους reference modes στα cluster maps.

\section{Χάρτες Συσταδοποίησης στο Επίπεδο $(\sigma,f)$ $-$ Cluster Maps}

Τα cluster maps απεικονίζουν τις screened εκτιμήσεις και τις συστάδες που
προκύπτουν από κάθε μέθοδο στο επίπεδο $(\sigma,f)$. Για λόγους
αναγνωσιμότητας, παρουσιάζεται η αντιπροσωπευτική περίπτωση του πρώτου
order sweep στην Περιοχή 3 στα Σχήματα~\ref{fig:amb:cluster-maps-area3-partitioning} έως \ref{fig:amb:cluster-maps-area3-agglomerative}, η οποία περιλαμβάνει έξι reference modes. Τα υπόλοιπα selected maps παρουσιάζονται αντίστοιχα για όλους τους συνδυασμούς
sweep και περιοχής ελέγχου στο Παράρτημα~\ref{app:amb:cluster-maps}.

\newcommand{\mainclusterfigure}[3]{%
\begin{figure}[p]
    \centering
    \includegraphics[width=0.94\textwidth,height=0.82\textheight,keepaspectratio]{#1}
    \caption{Selected cluster map για το πρώτο order sweep και την Περιοχή 3: #2.}
    #3
\end{figure}
}

\mainclusterfigure{\ambientfigurepath Ambient_Mag0.1_T600s_dt10ms_seed1997/orders1/clustering/by_control_area/area_3/kmeans/pdf/k-means_selected_cluster_map.pdf}{\textit{k-Means}}{\label{fig:amb:cluster-maps-area3-partitioning}}
\mainclusterfigure{\ambientfigurepath Ambient_Mag0.1_T600s_dt10ms_seed1997/orders1/clustering/by_control_area/area_3/kmedoids/pdf/k-medoids_selected_cluster_map.pdf}{\textit{k-Medoids}}{}
\mainclusterfigure{\ambientfigurepath Ambient_Mag0.1_T600s_dt10ms_seed1997/orders1/clustering/by_control_area/area_3/hdbscan/pdf/hdbscan_selected_cluster_map.pdf}{HDBSCAN}{}
\mainclusterfigure{\ambientfigurepath Ambient_Mag0.1_T600s_dt10ms_seed1997/orders1/clustering/by_control_area/area_3/gmm/pdf/gmm_selected_cluster_map.pdf}{GMM}{}
\mainclusterfigure{\ambientfigurepath Ambient_Mag0.1_T600s_dt10ms_seed1997/orders1/clustering/by_control_area/area_3/dbscan/pdf/dbscan_selected_cluster_map.pdf}{DBSCAN}{\label{fig:amb:cluster-maps-area3-density}}
\mainclusterfigure{\ambientfigurepath Ambient_Mag0.1_T600s_dt10ms_seed1997/orders1/clustering/by_control_area/area_3/optics/pdf/optics_selected_cluster_map.pdf}{OPTICS}{}
\mainclusterfigure{\ambientfigurepath Ambient_Mag0.1_T600s_dt10ms_seed1997/orders1/clustering/by_control_area/area_3/agglomerative/pdf/agglomerative_selected_cluster_map.pdf}{Agglomerative Hierarchical Clustering}{\label{fig:amb:cluster-maps-area3-agglomerative}}
\clearpage

Στα Σχήματα~\ref{fig:amb:cluster-maps-area3-partitioning} έως
\ref{fig:amb:cluster-maps-area3-agglomerative} φαίνεται η διαφορετική πολιτική
ανάθεσης των εκτιμήσεων. Οι \textit{k-Means} και \textit{k-Medoids} αποδίδουν
το σύνολο των σημείων σε τέσσερις συστάδες. Ο \textit{OPTICS} σχηματίζει επίσης
τέσσερις συστάδες, χαρακτηρίζοντας όμως το 62.9\% των εκτιμήσεων ως θόρυβο,
ενώ ο \textit{HDBSCAN} διατηρεί πέντε συστάδες. Ο \textit{GMM}, μετά το
post-processing της Ενότητας~\ref{amb:silhouette-filtering}, οδηγεί σε έξι
συστάδες, ίσες δηλαδή με το πλήθος των reference modes. Αντίθετα, ο
\textit{DBSCAN} διατηρεί μόνο δύο συστάδες, χαρακτηρίζοντας όμως πολύ μεγάλο
μέρος των εκτιμήσεων ως θόρυβο. Τέλος, η Agglomerative μέθοδος, όπως και ο
\textit{GMM}, διατηρεί έξι συστάδες.

Παρά τις διαφορές αυτές, οι θέσεις αρκετών συστάδων συμφωνούν ικανοποιητικά
ως προς τη συχνότητα με τα αντίστοιχα reference modes. Η συμφωνία αυτή
παρατηρείται σε περισσότερες από μία μεθόδους και δείχνει ότι οι κυρίαρχες
συχνότητες ορισμένων ρυθμών ανακτώνται σταθερά από τις screened εκτιμήσεις.
Ταυτόχρονα, σχεδόν όλες οι μέθοδοι σχηματίζουν μία ή περισσότερες συστάδες
σε πολύ χαμηλές συχνότητες, όπου δεν αντιστοιχεί reference mode, καθώς και
συστάδες σε υψηλές συχνότητες επίσης χωρίς αντίστοιχο mode αναφοράς. Οι
συστάδες αυτές υποδηλώνουν επαναλαμβανόμενες εκτιμήσεις μη
αντιπροσωπευτικών ρυθμών. Οι μέθοδοι που επιτρέπουν θόρυβο, δηλαδή DBSCAN, OPTICS και HDBSCAN, απορρίπτουν
πολύ μεγάλο μέρος των εκτιμήσεων ως θόρυβο \textendash{} ιδιαίτερα οι DBSCAN
και OPTICS. Επομένως, η καθαρότερη γεωμετρική εικόνα τους και οι υψηλές
τιμές silhouette πρέπει να συνεκτιμώνται με τη μειωμένη κάλυψη των
διαθέσιμων εκτιμήσεων.

\section{Ποσοτική Αξιολόγηση των Αποτελεσμάτων Συσταδοποίησης}

Η ποσοτική αξιολόγηση συμπληρώνει την οπτική σύγκριση των cluster maps.
Ο αριθμός συστάδων παρουσιάστηκε ήδη στον Πίνακα~\ref{tab:amb:clustering-comparison}.
Στη συνέχεια, η αξιολόγηση οργανώνεται σε τρία επίπεδα: τη συνοχή και την κάλυψη των εκτιμήσεων, τη συμφωνία με τα reference modes
και την ακρίβεια της θέσης των εκτιμήσεων που αντιστοιχίζονται σε αυτούς.

\subsection{Μετρικές Συνοχής και Κάλυψης}

Ο Πίνακας~\ref{tab:amb:cohesion-coverage} παρουσιάζει το τελικό πλήθος συστάδων, τον αριθμό των σημείων θορύβου, το ποσοστό των εκτιμήσεων που αποδίδονται σε συστάδα και το τελικό \textit{silhouette score}. Οι \textit{k-Means}, \textit{k-Medoids} και \textit{AHC} έχουν υψηλά ποσοστά αποδιδόμενων εκτιμήσεων, με χαμηλότερες όμως τιμές silhouette σε σχέση με τις περισσότερες density-based μεθόδους. Στον GMM, η τελική κάλυψη μεταβάλλεται λόγω της εφαρμογής του φίλτρου silhouette, ενώ οι DBSCAN και OPTICS εμφανίζουν γενικά υψηλές τιμές silhouette, αλλά το ποσοστό κάλυψης μεταβάλλεται σημαντικά μεταξύ των περιοχών και των order sweeps. Ο HDBSCAN διατηρεί γενικά μεγαλύτερο μέρος του δείγματος από τις άλλες density-based μεθόδους, με υψηλές τιμές silhouette.

\begingroup
\setlength{\tabcolsep}{4pt}
\setlength{\LTleft}{\fill}
\setlength{\LTright}{\fill}
\begin{longtable}{cccccccc}
\caption{Μετρικές συνοχής και κάλυψης της επιλεγμένης λύσης.}\label{tab:amb:cohesion-coverage}\\
\toprule
Sweep & Περ. & \makecell{Ref.\\modes} & Μέθοδος & Clusters & \makecell{Αποκλει-\\σμένα} & \makecell{Assigned\\(\%)} & Silhouette \\
\midrule
\endfirsthead
\toprule
Sweep & Περ. & \makecell{Ref.\\modes} & Μέθοδος & Clusters & \makecell{Αποκλει-\\σμένα} & \makecell{Assigned\\(\%)} & Silhouette \\
\midrule
\endhead
\endfoot
\bottomrule
\endlastfoot
\multirow{7}{*}{orders1} & \multirow{7}{*}{1} & \multirow{7}{*}{5} & k-Means & 5 & 0 & 100.0 & 0.640 \\*
& & & k-Medoids & 5 & 0 & 100.0 & 0.634 \\*
& & & DBSCAN & 2 & 335 & 40.7 & 0.876 \\*
& & & OPTICS & 2 & 492 & 12.9 & 0.952 \\*
& & & HDBSCAN & 4 & 165 & 70.8 & 0.796 \\*
& & & GMM & 7 & 143 & 74.7 & 0.621 \\*
& & & AHC & 5 & 2 & 99.6 & 0.633 \\
\midrule
\multirow{7}{*}{orders1} & \multirow{7}{*}{2} & \multirow{7}{*}{3} & k-Means & 4 & 0 & 100.0 & 0.700 \\*
& & & k-Medoids & 4 & 0 & 100.0 & 0.702 \\*
& & & DBSCAN & 3 & 140 & 50.0 & 0.927 \\*
& & & OPTICS & 5 & 124 & 55.7 & 0.922 \\*
& & & HDBSCAN & 7 & 60 & 78.6 & 0.782 \\*
& & & GMM & 7 & 70 & 75.0 & 0.734 \\*
& & & AHC & 4 & 0 & 100.0 & 0.685 \\
\midrule
\multirow{7}{*}{orders1} & \multirow{7}{*}{3} & \multirow{7}{*}{6} & k-Means & 4 & 0 & 100.0 & 0.650 \\*
& & & k-Medoids & 4 & 0 & 100.0 & 0.654 \\*
& & & DBSCAN & 2 & 389 & 29.5 & 0.936 \\*
& & & OPTICS & 4 & 347 & 37.1 & 0.935 \\*
& & & HDBSCAN & 5 & 140 & 74.6 & 0.816 \\*
& & & GMM & 6 & 165 & 70.1 & 0.718 \\*
& & & AHC & 6 & 0 & 100.0 & 0.633 \\
\midrule
\multirow{7}{*}{orders2} & \multirow{7}{*}{1} & \multirow{7}{*}{5} & k-Means & 2 & 0 & 100.0 & 0.609 \\*
& & & k-Medoids & 2 & 0 & 100.0 & 0.609 \\*
& & & DBSCAN & 5 & 214 & 60.9 & 0.716 \\*
& & & OPTICS & 6 & 349 & 36.3 & 0.857 \\*
& & & HDBSCAN & 9 & 187 & 65.9 & 0.741 \\*
& & & GMM & 6 & 317 & 42.2 & 0.705 \\*
& & & AHC & 2 & 1 & 99.8 & 0.587 \\
\midrule
\multirow{7}{*}{orders2} & \multirow{7}{*}{2} & \multirow{7}{*}{3} & k-Means & 5 & 0 & 100.0 & 0.681 \\*
& & & k-Medoids & 4 & 0 & 100.0 & 0.664 \\*
& & & DBSCAN & 4 & 58 & 78.7 & 0.774 \\*
& & & OPTICS & 7 & 113 & 58.5 & 0.817 \\*
& & & HDBSCAN & 4 & 34 & 87.5 & 0.749 \\*
& & & GMM & 8 & 27 & 90.1 & 0.587 \\*
& & & AHC & 4 & 0 & 100.0 & 0.670 \\
\midrule
\multirow{7}{*}{orders2} & \multirow{7}{*}{3} & \multirow{7}{*}{6} & k-Means & 2 & 0 & 100.0 & 0.621 \\*
& & & k-Medoids & 4 & 0 & 100.0 & 0.621 \\*
& & & DBSCAN & 5 & 201 & 64.8 & 0.752 \\*
& & & OPTICS & 4 & 422 & 26.1 & 0.952 \\*
& & & HDBSCAN & 6 & 224 & 60.8 & 0.742 \\*
& & & GMM & 4 & 378 & 33.8 & 0.821 \\*
& & & AHC & 6 & 0 & 100.0 & 0.607 \\
\end{longtable}
\endgroup

\subsection{Συμφωνία με τους Ρυθμούς Αναφοράς}
Για τον υπολογισμό των δεικτών \textit{V-measure} και ARI, απαιτείται η ύπαρξη ενός \textit{reference label} για κάθε εκτίμηση που δεν έχει χαρακτηριστεί ως θόρυβος. Αυτό ορίζεται ως το reference mode του \textit{PowerFactory} που βρίσκεται πιο κοντά στην εκτίμηση που δεν χαρακτηρίστηκε ως θόρυβος στο επίπεδο $(\sigma,\omega)$, με βάση την ευκλείδια απόσταση. Στη συνέχεια, υπολογίζονται οι δείκτες και παρουσιάζονται στον Πίνακα~\ref{tab:amb:v-measure_and_ari}. Και για τους δύο δείκτες, υψηλότερες τιμές υποδηλώνουν καλύτερη αντιστοίχηση με τα reference modes. Οι τιμές V-measure κυμαίνονται από $0$ έως $1$, ενώ οι τιμές ARI μπορούν να είναι και αρνητικές, με το $1$ να υποδηλώνει πλήρη συμφωνία και το $0$ να υποδηλώνει συμφωνία που δεν είναι καλύτερη από την τελείως τυχαία κατανομή.


\begin{longtable}{cccc}
\caption{Συμφωνία της επιλεγμένης συσταδοποίησης με τα τοπικά reference labels.}\label{tab:amb:v-measure_and_ari}\\
\toprule
Sweep & Μέθοδος & V-measure & ARI \\
\midrule
\endfirsthead
\toprule
Sweep & Μέθοδος & V-measure & ARI \\
\midrule
\endhead
\endfoot
\bottomrule
\endlastfoot
orders1\textendash{}Περ. 1 & k-Means & 0.685 & 0.515 \\
& k-Medoids & 0.689 & 0.530 \\
& DBSCAN & 0.856 & 0.894 \\
& OPTICS & 0.906 & 0.945 \\
& HDBSCAN & 0.784 & 0.652 \\
& GMM & 0.714 & 0.433 \\
& AHC & 0.682 & 0.519 \\
\midrule
orders1\textendash{}Περ. 2 & k-Means & 0.422 & 0.258 \\
& k-Medoids & 0.419 & 0.253 \\
& DBSCAN & 0.615 & 0.379 \\
& OPTICS & 0.573 & 0.354 \\
& HDBSCAN & 0.490 & 0.288 \\
& GMM & 0.491 & 0.295 \\
& AHC & 0.430 & 0.270 \\
\midrule
orders1\textendash{}Περ. 3 & k-Means & 0.715 & 0.621 \\
& k-Medoids & 0.739 & 0.642 \\
& DBSCAN & 1.000 & 1.000 \\
& OPTICS & 0.853 & 0.763 \\
& HDBSCAN & 0.778 & 0.649 \\
& GMM & 0.736 & 0.555 \\
& AHC & 0.693 & 0.603 \\
\midrule
orders2\textendash{}Περ. 1 & k-Means & 0.654 & 0.639 \\
& k-Medoids & 0.654 & 0.639 \\
& DBSCAN & 0.809 & 0.680 \\
& OPTICS & 0.619 & 0.309 \\
& HDBSCAN & 0.626 & 0.315 \\
& GMM & 0.601 & 0.285 \\
& AHC & 0.556 & 0.442 \\
\midrule
orders2\textendash{}Περ. 2 & k-Means & 0.409 & 0.277 \\
& k-Medoids & 0.420 & 0.275 \\
& DBSCAN & 0.374 & 0.221 \\
& OPTICS & 0.451 & 0.185 \\
& HDBSCAN & 0.424 & 0.278 \\
& GMM & 0.417 & 0.206 \\
& AHC & 0.449 & 0.371 \\
\midrule
orders2\textendash{}Περ. 3 & k-Means & 0.590 & 0.567 \\
& k-Medoids & 0.760 & 0.683 \\
& DBSCAN & 0.871 & 0.918 \\
& OPTICS & 0.858 & 0.727 \\
& HDBSCAN & 0.726 & 0.532 \\
& GMM & 0.889 & 0.839 \\
& AHC & 0.743 & 0.677 \\
\end{longtable}

Οι μέθοδοι k-Means και k-Medoids εμφανίζουν τιμές
\mbox{\textit{V-measure}} και ARI από μέτριες έως υψηλές, επομένως
εντοπίζουν σε σημαντικό βαθμό τα reference modes. Ο \mbox{HDBSCAN} παρουσιάζει παρόμοια εικόνα,
με καλή συμφωνία σε αρκετές περιοχές και ταυτόχρονα μεγαλύτερη κάλυψη
εκτιμήσεων από τις άλλες density-based μεθόδους. Ο GMM εμφανίζει επίσης
μεταβλητή συμφωνία με χαρακτηριστικό παράδειγμα τις υψηλές τιμές του στο δεύτερο sweep της Περιοχής~3, οι οποίες 
πρέπει να ερμηνεύονται μαζί με το μικρό ποσοστό εκτιμήσεων που διατηρείται
μετά το φίλτρο silhouette. Οι DBSCAN και OPTICS παρουσιάζουν μεταβλητή
συμφωνία με τα reference labels, ανάλογα με το order sweep και την περιοχή
ελέγχου. Η υψηλότερη συμφωνία εμφανίζεται στο πρώτο order sweep, ενώ στο
δεύτερο παρατηρούνται χαμηλότερες τιμές σε ορισμένες περιοχές. Πλήρης
συμφωνία επιτυγχάνεται μόνο από τον DBSCAN στο πρώτο order sweep της
Περιοχής~3. Οι δείκτες υπολογίζονται αποκλειστικά για τις εκτιμήσεις που
δεν χαρακτηρίζονται ως θόρυβος και, επομένως, πρέπει να ερμηνεύονται μαζί
με το ποσοστό κάλυψης και το πλήθος των συστάδων.

\subsection{Ακρίβεια Θέσης των Συστάδων\textendash{}Δείκτης MAD}

Για κάθε εκτίμηση που έχει αποδοθεί σε συστάδα υπολογίζεται η απόστασή της
από τον reference mode στον οποίο αντιστοιχίζεται. Ως δείκτης ακρίβειας
χρησιμοποιείται η διάμεση απόσταση (MAD) στο επίπεδο $(\sigma,\omega)$.
Μικρότερες τιμές υποδηλώνουν ότι οι εκτιμήσεις βρίσκονται εγγύτερα στον
αντίστοιχο reference mode. Οι Πίνακες~\ref{tab:amb:mad-orders1} και
\ref{tab:amb:mad-orders2} παρουσιάζουν τις αναλυτικές τιμές ανά reference
mode και μέθοδο. Κάθε κελί έχει τη μορφή \textit{πλήθος εκτιμήσεων / MAD}, με το MAD
σε rad/s. Η μεγάλη παύλα δηλώνει ότι δεν αντιστοιχίσθηκε καμία εκτίμηση
στον συγκεκριμένο reference mode.

\begin{table}[H]
    \centering
    \caption{Αναλυτικές τιμές MAD για το πρώτο order sweep (orders1).}
    \label{tab:amb:mad-orders1}
    \resizebox{\textwidth}{!}{%
    \begin{tabular}{cccccccccc}
        \toprule
        Μέθοδος & Mode 1 & Mode 2 & Mode 3 & Mode 4 & Mode 5 & Mode 6 & Mode 7 & Mode 8 & Mode 9 \\
        \midrule
        k-Means & \textemdash{} & 515 / 0.492 & 523 / 0.580 & \textemdash{} & 65 / 3.431 & \textemdash{} & \textemdash{} & 147 / 0.728 & 147 / 1.442 \\
        k-Medoids & \textemdash{} & 494 / 0.459 & 532 / 0.585 & \textemdash{} & 65 / 3.431 & \textemdash{} & \textemdash{} & 148 / 0.739 & 158 / 1.342 \\
        DBSCAN & \textemdash{} & 269 / 0.135 & 240 / 0.326 & \textemdash{} & 24 / 4.018 & \textemdash{} & \textemdash{} & \textemdash{} & \textemdash{} \\
        OPTICS & \textemdash{} & 184 / 0.100 & 134 / 0.530 & 29 / 0.888 & 49 / 1.051 & \textemdash{} & \textemdash{} & \textemdash{} & 38 / 2.492 \\
        HDBSCAN & \textemdash{} & 414 / 0.332 & 385 / 0.411 & \textemdash{} & 65 / 3.261 & \textemdash{} & \textemdash{} & 74 / 0.507 & 94 / 1.251 \\
        GMM & \textemdash{} & 404 / 0.382 & 328 / 0.368 & 27 / 0.888 & 76 / 3.180 & \textemdash{} & \textemdash{} & 62 / 0.500 & 122 / 1.299 \\
        AHC & \textemdash{} & 516 / 0.492 & 333 / 0.608 & \textemdash{} & 65 / 3.431 & 192 / 0.915 & \textemdash{} & 146 / 0.722 & 143 / 1.442 \\
        \bottomrule
    \end{tabular}
    }
\end{table}

\begin{table}[H]
    \centering
    \caption{Αναλυτικές τιμές MAD για το δεύτερο order sweep (orders2).}
    \label{tab:amb:mad-orders2}
    \resizebox{\textwidth}{!}{%
    \begin{tabular}{cccccccccc}
        \toprule
        Μέθοδος & Mode 1 & Mode 2 & Mode 3 & Mode 4 & Mode 5 & Mode 6 & Mode 7 & Mode 8 & Mode 9 \\
        \midrule
        k-Means & \textemdash{} & 606 / 0.832 & 216 / 2.515 & \textemdash{} & 56 / 3.749 & \textemdash{} & 290 / 1.939 & 223 / 1.730 & \textemdash{} \\
        k-Medoids & \textemdash{} & 590 / 0.815 & 361 / 1.190 & \textemdash{} & 54 / 3.720 & \textemdash{} & \textemdash{} & 223 / 1.730 & 163 / 1.342 \\
        DBSCAN & \textemdash{} & 406 / 0.445 & 311 / 0.468 & 34 / 0.712 & 32 / 4.044 & \textemdash{} & \textemdash{} & 38 / 0.641 & 97 / 1.336 \\
        OPTICS & \textemdash{} & 228 / 0.420 & 152 / 2.567 & 32 / 0.707 & 39 / 0.561 & \textemdash{} & \textemdash{} & 22 / 0.645 & 34 / 1.336 \\
        HDBSCAN & \textemdash{} & 342 / 0.597 & 329 / 0.494 & 48 / 0.764 & 49 / 3.749 & \textemdash{} & \textemdash{} & 74 / 0.729 & 104 / 1.228 \\
        GMM & \textemdash{} & 205 / 2.096 & 300 / 0.513 & \textemdash{} & 64 / 3.493 & \textemdash{} & \textemdash{} & 28 / 0.503 & 72 / 2.078 \\
        AHC & \textemdash{} & 677 / 0.919 & 219 / 2.542 & \textemdash{} & 53 / 3.768 & 165 / 0.963 & \textemdash{} & 133 / 1.301 & 143 / 1.373 \\
        \bottomrule
    \end{tabular}
    }
\end{table}
\clearpage

Οι τιμές MAD στους αναλυτικούς πίνακες είναι γενικά μη αμελητέες. Αυτό δείχνει ότι, ακόμη
και στις περιπτώσεις όπου σχηματίζεται συστάδα σε συχνότητα κοντά σε έναν
reference mode, παραμένει διασπορά ως προς τη συχνότητα ή/και την απόσβεση.
Η απόσταση υπολογίζεται στο επίπεδο $(\sigma,\omega)$ και συνεπώς δεν
εκφράζει μόνο απόκλιση στη συχνότητα. Επιπλέον, μικρότερη τιμή MAD δεν
αρκεί από μόνη της για την ανάδειξη μίας μεθόδου ως καταλληλότερης καθώς μπορεί
να συνδέεται με μικρό πλήθος διατηρούμενων εκτιμήσεων ή με πολύ μεγάλο αριθμό συστάδων, όπως στην περίπτωση του AHC. Επομένως, το MAD ερμηνεύεται μαζί
με το πλήθος των εκτιμήσεων κάθε κελιού, το ποσοστό αποδιδόμενων εκτιμήσεων και τους δείκτες V-measure και ARI.

\section{Συνολική Συζήτηση των Αποτελεσμάτων Συσταδοποίησης}

Η σύγκριση των μεθόδων δείχνει ότι δεν προκύπτει μία καθολικά βέλτιστη
λύση για όλους τους στόχους της συσταδοποίησης. Οι μέθοδοι
διαμερισμού, \textit{k-Means} και \textit{k-Medoids}, αποδίδουν το σύνολο
των screened εκτιμήσεων σε σχετικά μικρό αριθμό συστάδων και παρουσιάζουν
σταθερές, ενδιάμεσες τιμές συνοχής και συμφωνίας με τα reference modes.
Αποτελούν, επομένως, μία συνεπή επιλογή όταν προτεραιότητα είναι η
διατήρηση της πληροφορίας του διαθέσιμου δείγματος. Στον GMM, το φίλτρο
silhouette απορρίπτει τις συστάδες χαμηλής συνοχής, με αποτέλεσμα η τελική
κάλυψη να μεταβάλλεται ανάλογα με το order sweep και την περιοχή ελέγχου.
Ο AHC οδηγεί σε δύο έως έξι τελικές συστάδες και διατηρεί σχεδόν το σύνολο
των εκτιμήσεων. Επομένως, παρέχει λύση υψηλής κάλυψης με σχετικά μικρό
αριθμό τελικών συστάδων.

Οι density-based μέθοδοι παρουσιάζουν διαφορετικό συμβιβασμό μεταξύ
συνοχής και κάλυψης. Ο DBSCAN σχηματίζει από δύο έως πέντε συστάδες, ενώ ο
OPTICS από δύο έως επτά. Και οι δύο μέθοδοι εμφανίζουν γενικά υψηλές τιμές
silhouette και αναδεικνύουν συνεκτικές modal συγκεντρώσεις, χαρακτηρίζουν
όμως σημαντικό ποσοστό των εκτιμήσεων ως θόρυβο. Ο HDBSCAN διατηρεί
μεγαλύτερο μέρος του δείγματος και απορρίπτει παράλληλα τις εκτιμήσεις που
δεν ανήκουν σε πυκνές περιοχές, παρουσιάζοντας πιο ισορροπημένη
συμπεριφορά μεταξύ κάλυψης και συνοχής.

Η εξέταση των cluster maps και των τιμών MAD επιβεβαιώνει ότι αρκετοί
ρυθμοί αναφοράς ανακτώνται σε ικανοποιητική περιοχή συχνοτήτων. Παράλληλα,
παρατηρούνται συστάδες πολύ χαμηλών ή υψηλών συχνοτήτων χωρίς αντίστοιχο
ρυ\-θμό αναφοράς, καθώς και μη αμελητέα διασπορά ως προς τη συχνότητα ή την
απόσβεση. Επομένως, η επιλογή μεθόδου δεν μπορεί να βασιστεί σε μία μόνο
μετρική. Για εφαρμογές όπου απαιτείται ισορροπία μεταξύ διατήρησης των
εκτιμήσεων και απόρριψης μη συνεκτικών σημείων, ο HDBSCAN αποτελεί
ενδιαφέρουσα επιλογή. Όταν η πλήρης κάλυψη των διαθέσιμων εκτιμήσεων είναι
αυστηρή απαίτηση, οι \textit{k-Means} και \textit{k-Medoids} παρέχουν πιο
άμεση και ερμηνεύσιμη κατάτμηση. Οι DBSCAN και OPTICS μπορούν να
χρησιμοποιηθούν συμπληρωματικά, ως μηχανισμοί ανάδειξης των πιο πυκνών
modal συγκεντρώσεων.
